\documentclass[10pt]{article}
\usepackage[utf8]{inputenc}
\usepackage[T1]{fontenc}
\usepackage{mathtools} % Charge amsmath.
\usepackage{amssymb}
\usepackage{circuitikz}
\usepackage{fontspec}
\usepackage{array}
\usepackage{float}

\setmainfont{Times New Roman}

%
% NOTE: schematics generated using the online tool: 'https://www.circuit2tikz.tf.fau.de/designer'.
% Special thank to them!
%

\title{How to evaluate the speed of a bike}
\author{QuBi}
\date{\today}

\begin{document}

\maketitle

The usual tachometers you can find for your bike are sluggish. They take forever to display the instantaneous speed. Probably because of some averaging process.\\

There must be a smarter way of approximating the speed. In 2026, we can certainly do a better refresh rate than less than 1 fps.\\



\section*{Description of the problem}

In this article we describe an attempt at finding a "good" way of estimating the speed based on the common bike speed-o-meter setup: a magnetic sensor ticking on every spin of the wheel.

\begin{circuitikz}
	\draw (5, 4.75) -- (5, 5.25);
	\draw (7.75, 2) -- (8.25, 2);
	\node[circ] at (5, 2){};
	\draw (4, 5) -- (11, 5);
	\node[currarrow, xscale=1.2, yscale=1.2](N1) at (11, 5){} node[anchor=west] at (N1.east){$u$};
	\draw (8, 1) -- (8, 8);
	\node[currarrow, rotate=90, xscale=1.2, yscale=1.2](N2) at (8, 8){} node[anchor=south] at (N2.east){$p$};
	\node[circ] at (8, 5){};
	\draw[line width=2pt, -latex] (5, 4.25) -- (5, 5);
	\draw[line width=2pt, -latex] (8, 4.25) -- (8, 5);
	\node[shape=rectangle, minimum width=4.965cm, minimum height=0.965cm] at (8, 9){} node[anchor=center, align=center, text width=4.577cm, inner sep=6pt] at (8, 9){\footnotesize normalised distance};
	\node[shape=rectangle, minimum width=0.465cm, minimum height=0.715cm](N3) at (5.25, 4.375){} node[anchor=center] at (N3.text){$1$};
	\node[shape=rectangle, minimum width=0.465cm, minimum height=0.715cm](N4) at (8.25, 4.375){} node[anchor=center] at (N4.text){$2$};
	\node[shape=rectangle, minimum width=3.465cm, minimum height=0.965cm] at (13.5, 5){} node[anchor=center, align=center, text width=3.077cm, inner sep=6pt] at (13.5, 5){\footnotesize normalised speed};
\end{circuitikz}

Therefore the kernel function must satisfy:
\begin{itemize}	
	\item $p(-1) = -1$
	\item $p(0) = 0$
\end{itemize}

But that doesn't yield any prediction capacity and is too permissive.


Let's get more precise. The dynamic resistance (so the gain in the I/V curve) $R_C$ can be expressed explicitly as a function of $I_C$ therefore giving a more refined description of the transfer curve.

For a BJT:
\begin{equation}
\begin{split}
	a	& = \left. \frac{\partial I_C}{\partial V_{BE}} \right|_{V_{BE} = V_{BE}^0}\\
 		& = \left. \frac{\partial }{\partial V_{BE}} \right|_{V_{BE} = V_{BE}^0} I_S \exp \left( \right)\\
		& = \left( 1 - \dfrac{2R_C}{2R_C + Z} \right) \Delta I\\
		& = \dfrac{Z}{2R_C + Z} \Delta I
\end{split}
	a = 
\end{equation}

 
\subsection*{Limits}


\subsubsection*{Large signal analysis}
Unlike the Emitter Follower, there's no feedback here to rely on. Hence, we expect deviation from the model if the input is too wide.

\subsubsection*{Transistor spec mismatch}
This section is TODO.



\section*{Example 12: the diode Ladder Filter}

\subsubsection*{Description}
As a follow-up of the Moog Ladder Filter, we now consider the case of the diode version. It looks very similar to the Moog one, except the BJTs have been replaced with diodes. 



\section*{Example 13: the Moog Voltage Controlled High-Pass filter (module \textit{904B})}

\subsubsection*{Description}
Similarly, Moog came up with a design for a high-pass filter (also patented). The structure differs a bit from the low-pass, but it is equally as puzzling.

Here below is one stage of the high-pass filter:






\end{document}